% STATUT : COMPLÉTER PLUS TARD
% SOURCE : PROJECT_TRUTH.md ; guide EPI
% LIMITE : 2 pages maximum dans la version finale.
\chapter*{Introduction générale}
\markboth{Introduction générale}{}
\addcontentsline{toc}{chapter}{Introduction générale}
\label{chap:00-introduction}
Le stage porte sur la conception et la réalisation d'une application web interne destinée à organiser les demandes techniques reçues des clients. L'enjeu est de qualifier les demandes, attribuer les responsabilités, sécuriser les accès, suivre les transitions et conserver une trace des changements.

La problématique est la suivante : comment construire une solution simple à exploiter, maintenable et contrôlée pour accompagner tout le cycle de vie d'une demande, tout en conservant une décision humaine lorsque l'assistance IA est sollicitée ?

Le projet vise une architecture Angular, Spring Boot et PostgreSQL, structurée en modules et protégée par JWT et RBAC. Le rapport présente d'abord le contexte et le cadrage, puis l'analyse et la conception, avant de décrire la réalisation et la validation. Les informations administratives absentes restent signalées comme éléments à compléter.
